%==============================================================
%  lecons/ensembles/construction_Q.tex — Construction de Q
%==============================================================

\lecontitle{Construction de \texorpdfstring{$\mathbb{Q}$}{Q}}{Théorie des ensembles — Corps des fractions de \texorpdfstring{$\mathbb{Z}$}{Z}}

%=============================================================
%  § 1 — MOTIVATION
%=============================================================
\secnum{navyblue}{1}{Motivation — inverser la multiplication}

\rubriq{Pourquoi étendre $\mathbb{Z}$ ?}
Dans $\mathbb{Z}$, l'équation $n \cdot x = m$ n'a de solution entière que
si $n \mid m$. On veut le plus petit corps contenant $\mathbb{Z}$ — on force
la division par tout entier non nul. Ce corps est exactement
$\mathbb{Q} = \mathrm{Frac}(\mathbb{Z})$.

\begin{tcolorbox}[thmbox]
  \textbf{Définition.}\enspace%
  \index{nombres rationnels!construction}\index{corps des fractions}
  Sur $S = \mathbb{Z} \times \mathbb{Z}^*$ (avec $\mathbb{Z}^* = \mathbb{Z}\setminus\{0\}$),
  on pose
  \[
    (a,b) \sim (c,d) \;\iff\; ad = bc.
  \]
  $\mathbb{Q} = S/{\sim}$, muni des lois $\tfrac{a}{b}+\tfrac{c}{d}=\tfrac{ad+bc}{bd}$
  et $\tfrac{a}{b}\cdot\tfrac{c}{d}=\tfrac{ac}{bd}$, est un corps.
\end{tcolorbox}

\decsep{}

%=============================================================
%  § 2 — STRUCTURE ALGÉBRIQUE ET ORDRE
%=============================================================
\secnum{navyblue}{2}{Structure de $\mathbb{Q}$}

\begin{tcolorbox}[thmbox]
  \textbf{Propriétés de $\mathbb{Q}$.}\enspace%
  \index{nombres rationnels!propriétés}
  \begin{enumerate}
    \item $\mathbb{Q}$ est un \emph{corps commutatif} de caractéristique $0$.
    \item L'ordre $\tfrac{a}{b} \leq \tfrac{c}{d}$ (avec $b,d > 0$)
          $\Leftrightarrow ad \leq bc$ fait de $\mathbb{Q}$ un \emph{corps ordonné}.
    \item $\mathbb{Q}$ est \emph{dense} : entre deux rationnels distincts
          il y en a toujours un autre.
    \item $\mathbb{Q}$ est \emph{dénombrable} : il existe une bijection
          $\mathbb{N} \to \mathbb{Q}$ (argument de Cantor — diagonale).
  \end{enumerate}
\end{tcolorbox}

\begin{mdframed}[style=proofbar]
  \textit{Densité.} Si $p < q$ dans $\mathbb{Q}$, alors $p < \tfrac{p+q}{2} < q$. $\checkmark$

  \textit{Dénombrabilité.} On liste les fractions irréductibles $a/b$ avec $b>0$
  en parcourant les diagonales du tableau $\mathbb{Z} \times \mathbb{N}^*$.
  L'image de cette liste est $\mathbb{Q}$, donc $|\mathbb{Q}| \leq |\mathbb{N}|$.
  Comme $\mathbb{N} \hookrightarrow \mathbb{Q}$, on a $|\mathbb{Q}| = \aleph_0$.
  \hfill$\blacksquare$
\end{mdframed}

\rubriq{Représentation unique}
Toute fraction $a/b \in \mathbb{Q}$ s'écrit de façon unique sous la forme
$p/q$ avec $\pgcd(p,q) = 1$ et $q > 0$ (\emph{fraction irréductible}).

\decsep{}

%=============================================================
%  § 3 — INCOMPLÉTUDE DE Q
%=============================================================
\secnum{navyblue}{3}{Incomplétude de $\mathbb{Q}$}

\begin{tcolorbox}[thmbox]
  \textbf{Théorème ($\sqrt{2} \notin \mathbb{Q}$).}\enspace%
  \index{irrationalité de $\sqrt{2}$}
  Il n'existe pas de rationnel $x$ tel que $x^2 = 2$.
\end{tcolorbox}

\begin{mdframed}[style=proofbar]
  Supposons $x = p/q$ irréductible vérifiant $p^2 = 2q^2$.
  Alors $2 \mid p^2$, donc $2 \mid p$ (car $2$ est premier). Écrivons $p = 2k$.
  Alors $4k^2 = 2q^2$, soit $2k^2 = q^2$, donc $2 \mid q$. Contradiction avec
  $\pgcd(p,q) = 1$. \hfill$\blacksquare$
\end{mdframed}

\rubriq{Conséquence : $\mathbb{Q}$ n'est pas complet}
La suite $(x_n)$ définie par $x_0 = 1$ et $x_{n+1} = \tfrac{x_n}{2} + \tfrac{1}{x_n}$
est de Cauchy dans $\mathbb{Q}$ mais converge vers $\sqrt{2} \notin \mathbb{Q}$.
Cette incomplétude motive la construction de $\mathbb{R}$.

\decsep{}

%=============================================================
%  § 4 — EXEMPLES, PIÈGES ET EXERCICES
%=============================================================
\secnum{exgreen}{4}{Exemples, pièges et exercices}

\noindent{\bfseries\color{navyblue} Exemples.}

\begin{mdframed}[style=exbar]
  \textbf{1. Corps premier.}\enspace%
  \index{corps premier}
  $\mathbb{Q}$ est le \emph{corps premier} de caractéristique $0$ : tout corps
  de caractéristique $0$ contient un sous-corps isomorphe à $\mathbb{Q}$.

  \smallskip\noindent
  \textbf{2. Approximation.}\enspace%
  Théorème de Dirichlet : pour tout $\alpha \in \mathbb{R}$ et $N \geq 1$,
  il existe $p/q$ avec $1 \leq q \leq N$ et $|\alpha - p/q| \leq 1/(qN)$.
  Ainsi $\mathbb{Q}$ est dense dans $\mathbb{R}$ pour la topologie usuelle.

  \smallskip\noindent
  \textbf{3. Topologie.}\enspace%
  Muni de la valeur absolue usuelle, $\mathbb{Q}$ est un espace métrique.
  Sa complétion est $\mathbb{R}$. Muni de la valeur absolue $p$-adique
  $|\cdot|_p$, sa complétion est $\mathbb{Q}_p$.
\end{mdframed}

\vspace{6pt}
\noindent{\bfseries\color{counterred} Pièges.}

\begin{mdframed}[style=cexbar]
  \textbf{Dense $\neq$ complet.}\enspace%
  $\mathbb{Q}$ est dense dans $\mathbb{R}$ (tout réel est limite de rationnels),
  mais pas complet : des suites de Cauchy rationnelles peuvent converger vers
  des irrationnels.

  \smallskip\noindent
  \textbf{Corps ordonné $\neq$ archimédien.}\enspace%
  Tout corps ordonné n'est pas archimédien. Mais $\mathbb{Q}$ l'est :
  pour tout $x \in \mathbb{Q}^+$, il existe $n \in \mathbb{N}$ tel que $n > x$.
\end{mdframed}

\vspace{6pt}
\noindent{\bfseries\color{goldaccent} Exercices.}

\begin{mdframed}[style=exobar]
  \begin{enumerate}
    \item Montrer que $\sqrt{3}$, $\sqrt{5}$ et $\sqrt{2}+\sqrt{3}$ sont irrationnels.

    \item Montrer que $\mathbb{Q}$ est dénombrable en utilisant le théorème
          de Schröder-Bernstein.

    \item Montrer que tout corps ordonné archimédien se plonge dans $\mathbb{R}$.

    \item Montrer que $\mathbb{Q}$ n'admet pas d'extension algébrique finie
          de degré $1$ autre que lui-même (i.e.\ tout polynôme irréductible
          sur $\mathbb{Q}$ de degré $\geq 2$ n'a pas de racine dans $\mathbb{Q}$).
          Donner un exemple de degré $2$.

    \item Soit $\alpha \in \mathbb{R}$ algébrique sur $\mathbb{Q}$.
          Montrer que $\mathbb{Q}(\alpha)$ est dénombrable.
  \end{enumerate}
\end{mdframed}

\leconfooter{Euclide (irrationnalité de $\sqrt{2}$, \textit{ca.} $-$300) — G.\ Cantor (dénombrabilité, 1874)}
